\section{Related Work}
\label{sec:related-work}

\paragraph{LLM reranking paradigms.}
LLM-based reranking is usually framed through pointwise, pairwise,
listwise, or setwise prompting. Pointwise rerankers score each
query--document pair independently, either through supervised
sequence-to-sequence relevance prediction or zero-shot likelihood-style
scoring~\citep{nogueira2020monot5,sachan2022upr,ma2024rankllama}.
Pairwise methods ask the LLM to compare two documents and then aggregate
local preferences into a ranking; Pairwise Ranking Prompting (PRP) shows
that this simplified comparison task can be effective for open-weight
models, while PRP-Graph further exploits label-generation probabilities
and graph aggregation to improve robustness and effectiveness
~\citep{qin2024prp,luo2024prpgraph}. Listwise methods instead prompt the
LLM to emit a permutation over a candidate block, as in RankGPT and its
open-weight successors~\citep{sun2023rankgpt,ma2023zeroshotlistwise,
pradeep2023rankvicuna,pradeep2023rankzephyr,zhang2025rankwithoutgpt}.
Setwise prompting occupies a middle ground: each call compares more than
two documents, but the output is a short label rather than a complete
permutation~\citep{zhuang2024setwise}. MaxContext belongs to this last
family: it keeps the short-label setwise interface, but changes the
granularity of the comparison from a small local window to the whole
current live pool.

\paragraph{Setwise prompting and setwise extensions.}
\citet{zhuang2024setwise} introduce setwise prompting by presenting a
small candidate window, asking the model to select the most relevant
document, and embedding that decision inside heapsort or bubblesort. We
build on the public \texttt{llm-rankers} implementation
~\citep{ielab_llmrankers}, but replace the fixed letter-label interface
with numeric labels so that all documents in the current live pool are
addressable. Recent work extends setwise reranking along complementary
axes. Setwise Insertion uses the first-stage ranking as prior knowledge,
reducing unnecessary comparisons by focusing on candidates likely to
improve the current order~\citep{podolak2025setwiseinsertion}. Rank-R1
keeps the setwise selection interface but trains the model with
reinforcement learning to reason before choosing the most relevant
candidate~\citep{zhuang2025rankr1}. MaxContext is orthogonal to both
directions: it does not warm-start an insertion procedure from BM25, and
it does not train a new reasoning model. Instead, it asks what setwise
selection looks like when the candidate window is expanded to the whole
live reranking pool.

\paragraph{Tournament, graph, and recursive aggregation methods.}
Several recent rerankers try to extract more global information from
multi-document or pairwise judgments by changing the comparison schedule
or aggregation rule. TourRank uses a tournament-inspired grouping
strategy and a points system to mitigate context-length limits, reduce
latency, and improve robustness to document input order
~\citep{chen2025tourrank}. BLITZRANK formalizes $k$-wise reranking as a
tournament-graph problem: each $k$-item comparison induces many pairwise
preferences, which can be aggregated through a global preference graph
and transitive closure~\citep{agrawal2026blitzrank}. PRP-Graph similarly
constructs and aggregates pairwise preference graphs, but starts from
pairwise PRP units and uses target-label probabilities to express
comparison confidence~\citep{luo2024prpgraph}. REALM models relevance as
uncertain Gaussian estimates and recursively updates them with Bayesian
style evidence aggregation~\citep{wang2025realm}. RefRank reduces the
cost of comparative ranking by evaluating each document against one or
more reference anchors rather than comparing all document pairs
~\citep{li2025refrank}. These methods are close in motivation to
MaxContext because they all revisit the cost--quality tradeoff of LLM
reranking. The distinction is structural: TourRank, BLITZRANK,
PRP-Graph, REALM, and RefRank improve how multiple judgments are
scheduled, weighted, or aggregated, whereas MaxContext changes the
setwise selector's live-pool granularity while preserving a direct
extreme-selection output.

\paragraph{Listwise reranking and long-context ordering.}
Listwise rerankers ask the model to generate an ordered list of document
identifiers~\citep{sun2023rankgpt,ma2023zeroshotlistwise,
pradeep2023rankvicuna,pradeep2023rankzephyr}. Because a full permutation
is a longer and more brittle output object than a single label, listwise
systems often rely on sliding windows, permutation distillation, or
specialized open-weight training~\citep{sun2023rankgpt,
zhang2025rankwithoutgpt}. Efficiency-oriented listwise work modifies the
decoding or scoring interface: FIRST ranks candidates using the logits
of the first generated identifier and adds a learning-to-rank objective
~\citep{reddy2024first}, while self-calibrated listwise reranking
estimates globally comparable relevance scores instead of directly
emitting only a textual permutation~\citep{ren2025selfcalibrated}.
MaxContext differs from these methods in both output type and claim
scope. It uses long context for input visibility, but still asks for one
or two labels rather than a full permutation. Its deterministic
call-count claim is therefore made against matched whole-pool
single-extreme setwise controls, not against listwise rerankers.

\paragraph{Steering and training-based LLM rankers.}
Another line of work improves LLM ranking by changing the model rather
than the comparison schedule. RankLLaMA and related systems fine-tune
decoder-only LLMs for pointwise reranking in multi-stage retrieval
pipelines~\citep{ma2024rankllama}. Rank-without-GPT and RankZephyr show
that open-weight listwise rerankers can be built through distillation or
instruction tuning rather than relying on proprietary GPT models
~\citep{zhang2025rankwithoutgpt,pradeep2023rankzephyr}. Rank-R1 trains a
setwise reranker with reinforcement learning so that the model produces
reasoning before selecting a candidate~\citep{zhuang2025rankr1}. More
recently, RankSteer proposes post-hoc activation steering for pointwise
LLM ranking, using steering directions for decision, evidence, and role
signals to calibrate relevance judgments without explicit cross-document
comparison~\citep{wang2026ranksteer}. MaxContext is intentionally more
minimal: it is a zero-shot prompting method, does not modify model
weights or hidden states, and tests whether current long-context
open-weight models can support a whole-pool setwise interface directly.

\paragraph{Robustness, position effects, and efficiency accounting.}
Whole-pool prompting does not imply order invariance. Long-context models
can under-use information in some prompt positions~\citep{liu2024lostmiddle},
and listwise ranking can exhibit strong positional sensitivity; permutation
self-consistency addresses this by repeatedly shuffling inputs and
aggregating the resulting rankings~\citep{tang2024foundmiddle}. More
general studies of LLM multiple-choice and judge-style prompting also
show that option labels and presentation order can influence
outputs~\citep{zheng2024mcqbias,shi2025positionbias}. For this reason,
we include an input-order stability protocol rather than assuming that
whole-pool visibility makes MaxContext robust. We also separate
efficiency axes. Prior work emphasizes that LLM reranker efficiency
cannot be summarized by a single proxy, because calls, input tokens,
output tokens, FLOPs, batching, and wall-clock time can move in different
directions~\citep{zhuang2024setwise,podolak2025setwiseinsertion,
peng2025flopsreranking}. We therefore report LLM calls, wall-clock,
prompt/completion/total tokens, parser recovery, and BM25 bypasses, while
keeping the central MaxContext claim narrowly focused on serial LLM calls
and observed wall-clock behavior. Statistical comparisons follow
standard IR evaluation practice using paired tests and multiplicity
control~\citep{smucker2007comparison,urbano2019significance}.